Thesis: Growth Observables, Dual-Gate Torque Release, and Multipole Geometry at \(u=1\)

The structure-growth parameter  
\[
S_8=\sigma_8\left(\frac{\Omega_m}{0.3}\right)^{0.5}
\]  
is not an independent free quantity. It is a derived functional of the linear matter power spectrum at the present epoch, which itself is fixed once the background expansion history \(H(z)\) and the associated linear growth factor \(D(z)\) are specified. Within the Dual-Gate / Cyclic-Multipole framework the expansion history is completely determined by the late-time torque release, and that torque release is the direct numerical evaluation of the multipole geometry at the present-day boundary \(u=1\). The present thesis unfolds this chain with full mathematical explicitness.

1. Definition of the Structure-Growth Parameter

The root-mean-square amplitude of linear matter fluctuations in spheres of comoving radius \(8\,h^{-1}\mathrm{Mpc}\) is
\[
\sigma_8^2=\int\frac{dk}{k}\,\frac{k^3P(k)}{2\pi^2}\,W^2(kR_8),
\]  
where \(P(k)\) is the present-day linear matter power spectrum and \(W\) is the Fourier transform of the top-hat window. The combination that appears in weak-lensing analyses is conventionally written
\[
S_8=\sigma_8\left(\frac{\Omega_m}{0.3}\right)^{0.5}.
\]  
Both \(\sigma_8\) and \(\Omega_m\) are therefore functionals of the same underlying cosmological evolution.

2. Dependence on the Expansion History

The linear growth factor \(D(a)\) satisfies the ordinary differential equation
\[
\frac{d^2D}{da^2}+\left(\frac{3}{a}+\frac{1}{H}\frac{dH}{da}\right)\frac{dD}{da}-\frac{3}{2}\frac{\Omega_m(a)}{a^2}D=0,
\]  
with initial conditions fixed deep in the matter era. The Hubble function \(H(a)\) enters both the friction term and the source term through
\[
\Omega_m(a)=\frac{\Omega_m H_0^2}{a^3H(a)^2}.
\]  
Consequently every growth observable—\(\sigma_8\), the redshift-space distortion parameter \(f\sigma_8\), and the weak-lensing convergence power spectrum—is uniquely determined once \(H(z)\) is given.

3. Dual-Gate Expansion History

In the framework the effective Hubble parameter is
\[
H_{\rm eff}(z)=H_{\rm base}(z)+\Delta H_{\rm torque}(z),
\]  
where the base term recovers a standard flat-\(\Lambda\)CDM expansion with \(H_0^{\rm base}\approx70\) and the torque term is
\[
\Delta H_{\rm torque}(z)=\bigl({\rm sh\,geom\,base}\cdot a\cdot\sin\delta\phi_{\rm torque}\bigr)\times\mathcal{S}(u)\times0.528.
\]  
The numerical coefficients are fixed:
\[
{\rm sh\,geom\,base}=3.170,\qquad\delta\phi_{\rm torque}=1.72113420759.
\]  
The suppression factor \(\mathcal{S}(u)\) combines the exponential damping \(e^{-z/5.8}\) and the screening \(e^{-\chi/4500}\). At late times the torque approaches its full amplitude, elevating the local expansion rate to
\[
H_0=H_0^{\rm base}+3.17=73.17\,\mathrm{km\,s^{-1}\,Mpc^{-1}}.
\]

4. Torque Release as Evaluation of Multipole Geometry at \(u=1\)

The scaled clock coordinate is defined by
\[
u=\frac{N}{N_{\rm today}}=\frac{N}{10^9}\in[0,1].
\]  
At the present-day boundary \(u=1\) the cyclic multipole engine opens the complete hierarchy
\[
\sum_{\ell=0}^{4}(2\ell+1)=25.
\]  
These twenty-five states constitute the full set of projectable visibility channels. Their high-symmetry nodal vertices are the geometric loci at which residual screening vanishes. The Local Geometric Torque Contribution is precisely the numerical evaluation of that fully opened multipole structure:
\[
\Delta H_{\rm torque}(u=1)=3.17\,\mathrm{km\,s^{-1}\,Mpc^{-1}}.
\]  
Thus the late-time modification to \(H(z)\) is not an independent phenomenological insertion; it is the macroscopic consequence of the multipole geometry evaluated at the terminal lattice tick.

5. Closure of the Logical Chain

Because \(H_{\rm eff}(z)\) is fixed by the torque release, the linear growth equation is fixed. Because the growth equation is fixed, the present-day power spectrum \(P(k)\) is fixed. Because \(P(k)\) is fixed, both \(\sigma_8\) and \(\Omega_m\) are fixed, and therefore
\[
S_8=\sigma_8\left(\frac{\Omega_m}{0.3}\right)^{0.5}
\]  
is completely determined. The identical multipole geometry that produces \(H_0=73.17\) at \(u=1\) thereby determines the theory prediction that is compared with the DES Year 6 and KiDS-Legacy constraints after marginalization over intrinsic alignments and baryonic feedback.

The expansion history is fixed by the Dual-Gate torque.  
The Dual-Gate torque is the evaluation of the multipole geometry at \(u=1\).  
The growth observables that define \(S_8\) are functionals of that expansion history.  
